\chapter{Показательные и логарифмические уравнения и неравенства}\label{chap:exam-exponential}\markright{Глава 23}
\section{Основание и направление сравнения}
В показательных и логарифмических задачах сначала фиксируют основание. Для произвольного действительного показателя используем положительное основание. Если \(a=1\), выражение \(1^x\) равно единице, но логарифма по основанию \(1\) не существует.
\begin{definition}\term{Показательное уравнение} (exponential equation) содержит неизвестное в показателе степени.\end{definition}
\begin{definition}\term{Логарифмическое уравнение} (logarithmic equation) содержит логарифм выражения с неизвестным; неизвестное может находиться и в основании.\end{definition}
\begin{theorem}[Сравнение при постоянном основании]
Пусть \(a>0,a\ne1\). Равенство \(a^u=a^v\) равносильно \(u=v\). При \(a>1\) знак сравнения показателей сохраняется, при \(0<a<1\) меняется. Для \(u,v>0\) то же правило действует при сравнении \(\log_a u\) и \(\log_a v\).
\end{theorem}
\begin{proof}При \(a>1\) и \(v>u\) имеем \(\frac{a^v}{a^u}=a^{v-u}>1\), поэтому показательная функция строго возрастает. При \(0<a<1\) запишем \(a=\frac1b\), где \(b>1\); тогда \(\frac{a^v}{a^u}=b^{-(v-u)}<1\). Строгая монотонность даёт равносильность равенств. Логарифм является обратной функцией: если бы направление сравнения его значений было иным, применение показательной функции противоречило бы сравнению аргументов.\end{proof}
\par\addvspace{2pt}\begin{lrbox}{\CatFlowBox}\begin{minipage}[t]{\linewidth}
\setlength{\parskip}{2pt}\setlength{\columnsep}{4.5mm}\setlength{\intextsep}{0pt}
\begin{wrapfigure}[9]{r}{54mm}
\includegraphics[width=54mm,height=58mm,keepaspectratio]{images/unique/art-224.png}
\label{catflow:224}
\end{wrapfigure}
Направление неравенства определяется основанием, а не знаком показателя. Отрицательные значения показателя вполне допустимы: они дают положительные дробные степени основания. Перед сравнением полезно привести обе части к одному основанию и только после этого убрать степени. Если одинаковое основание получить неудобно, обе положительные части можно прологарифмировать по основанию больше единицы. При этом логарифмировать сумму по слагаемым нельзя. Выражение с несколькими степенями сначала пытаются упростить с помощью общего множителя или замены. Такой порядок помогает увидеть структуру задачи до вычислений и сохраняет все условия, от которых зависит направление перехода. В ответе остаётся проверить, не было ли дополнительных ограничений в исходных выражениях.
\par\WFclear\end{minipage}\end{lrbox}%
\ifdim\dimexpr\ht\CatFlowBox+\dp\CatFlowBox\relax<58mm\dp\CatFlowBox=\dimexpr58mm-\ht\CatFlowBox\relax\fi%
\noindent\usebox{\CatFlowBox}\par\addvspace{2pt}
\begin{example}Решим \((\frac13)^{x^2-2x}\geq\frac1{27}\). Справа стоит \((\frac13)^3\). Основание меньше единицы, поэтому \(x^2-2x\leq3\), то есть \((x-3)(x+1)\leq0\). Ответ \([-1;3]\).\end{example}

\section{Замены в показательных задачах}
\begin{definition}\term{Замена неизвестного} (substitution) вводит новое выражение \(t=\varphi(x)\), упрощающее уравнение. Вместе с заменой указывают допустимые значения \(t\) и способ возвращения к \(x\).\end{definition}
Для \(t=a^x\) при \(a>0,a\ne1\) имеем \(t>0\), и каждое положительное \(t\) даёт ровно одно \(x=\log_a t\). Например, \(a^{2x}=t^2\), \(a^{-x}=\frac1t\). При замене \(t=a^{x^2}\) положительности недостаточно: если \(a>1\), то \(t\geq1\), а при \(0<a<1\) имеем \(0<t\leq1\). Обратная подстановка может дать два, одно или ни одного значения \(x\).
\begin{example}Для \(9^x-10\cdot3^x+9=0\) положим \(t=3^x>0\). Получаем \((t-1)(t-9)=0\). Обе величины положительны; \(3^x=1\) даёт \(x=0\), а \(3^x=9\) даёт \(x=2\).\end{example}
\begin{example}В неравенстве \(4^x-6\cdot2^x+8\leq0\) замена \(t=2^x>0\) приводит к \((t-2)(t-4)\leq0\), то есть \(2\leq t\leq4\). Поскольку основание \(2\) больше единицы, возвращаемся к \(1\leq x\leq2\).\end{example}

\section{ОДЗ и преобразования логарифмов}
Логарифм \(\log_a F\) требует \(a>0,a\ne1,F>0\). Все эти ограничения накладывают до преобразований. Сумма \(\log_a F+\log_a G\) определена лишь при \(F>0,G>0\), хотя \(\log_a(FG)\) может иметь смысл и при двух отрицательных множителях. Поэтому объединять сумму можно с сохранением исходной ОДЗ. Наоборот, разложение логарифма произведения в сумму допустимо лишь там, где каждый множитель положителен.
\begin{theorem}[Квадрат под логарифмом]
При \(a>0,a\ne1,u\ne0\)
\[
\log_a(u^2)=2\log_a|u|.
\]
Без модуля правая часть определена только при \(u>0\).
\end{theorem}
\par\addvspace{2pt}\begin{lrbox}{\CatFlowBox}\begin{minipage}[t]{\linewidth}
\setlength{\parskip}{2pt}\setlength{\columnsep}{4.5mm}\setlength{\intextsep}{0pt}
\begin{wrapfigure}[11]{l}{54mm}
\includegraphics[width=54mm,height=58mm,keepaspectratio]{images/unique/art-227.png}
\label{catflow:227}
\end{wrapfigure}
\textit{Доказательство.} Число \(|u|\) положительно, а \(u^2=|u|^2\). Поэтому правило логарифма степени с положительным основанием аргумента даёт равенство.\hfill\(\blacksquare\)
\par\WFclear\end{minipage}\end{lrbox}%
\ifdim\dimexpr\ht\CatFlowBox+\dp\CatFlowBox\relax<58mm\dp\CatFlowBox=\dimexpr58mm-\ht\CatFlowBox\relax\fi%
\noindent\usebox{\CatFlowBox}\par\addvspace{2pt}

\par\addvspace{2pt}\begin{lrbox}{\CatFlowBox}\begin{minipage}[t]{\linewidth}
\setlength{\parskip}{2pt}\setlength{\columnsep}{4.5mm}\setlength{\intextsep}{0pt}
\begin{wrapfigure}[11]{l}{54mm}
\includegraphics[width=54mm,height=58mm,keepaspectratio]{images/unique/art-225.png}
\label{catflow:225}
\end{wrapfigure}
Запреты логарифма относятся ко всей исходной записи. Положительность аргумента одного логарифма не заменяет условий для остальных, а ненулевой знаменатель нельзя забывать после сокращения. Особенно внимательно проверяют вложенные логарифмы: внутренний логарифм должен не просто существовать, но и давать положительное число, если он служит аргументом внешнего. Замена через логарифм удобна тем, что переводит положительные значения исходной переменной во всю числовую прямую. Но после решения нового уравнения нужно вернуться к исходному неизвестному и проверить каждый знаменатель и каждое основание. Эта проверка коротка, если ограничения были выписаны заранее, и значительно труднее, если они восстанавливаются по уже упрощённой формуле.
\par\WFclear\end{minipage}\end{lrbox}%
\ifdim\dimexpr\ht\CatFlowBox+\dp\CatFlowBox\relax<58mm\dp\CatFlowBox=\dimexpr58mm-\ht\CatFlowBox\relax\fi%
\noindent\usebox{\CatFlowBox}\par\addvspace{2pt}
\begin{example}Решим \(\log_2(x-1)+\log_2(x+1)=3\). ОДЗ: \(x>1\). На ней объединение даёт \(\log_2(x^2-1)=3\), откуда \(x^2=9\). Из кандидатов \(-3,3\) подходит лишь \(3\). Если бы исходно стоял один логарифм произведения, отрицательный кандидат тоже мог бы оказаться допустимым; это другая область определения.\end{example}
При замене \(t=\log_a x\) допустимы все действительные \(t\), а обратная подстановка \(x=a^t\) автоматически даёт \(x>0\). Так \((\log_2x)^2-3\log_2x+2=0\) приводит к \(t=1,2\) и \(x=2,4\).

\section{Переменное основание и рационализация}
\begin{definition}\term{Рационализация неравенства} (sign-preserving rationalization) заменяет выражение другим с тем же знаком на заранее указанной ОДЗ. Это утверждение о знаках, а не о равенстве значений выражений.\end{definition}
\begin{theorem}[Точное правило рационализации]
В каждой точке, где \(a>0,a\ne1,F>0,G>0\), выражения
\[
\log_a F-\log_a G\quad\text{и}\quad(a-1)(F-G)
\]
имеют одинаковые знаки, включая нуль. Здесь \(a,F,G\) могут зависеть от \(x\). При \(a>0,a\ne1\) одинаковые знаки имеют также
\[
a^U-a^V\quad\text{и}\quad(a-1)(U-V).
\]
\end{theorem}
\begin{proof}Если \(a>1\), логарифм и степень строго возрастают, поэтому знак разности равен знаку разности аргументов или показателей; множитель \(a-1\) положителен. Если \(0<a<1\), обе функции строго убывают, а \(a-1<0\), что меняет оба знака. При равных аргументах либо показателях обе разности равны нулю.\end{proof}
В частности, \(\log_a F>c\) сначала записывают как \(\log_a F-\log_a(a^c)>0\), а затем получают \((a-1)(F-a^c)>0\). Все исходные условия на \(a,F\) сохраняются. Аналогично работают знаки \(<,\leq,\geq\). В произведении или частном допустимо заменять отдельные множители выражениями того же знака; у знаменателя обязательно сохраняется запрет нуля. Сумма не определяется знаками слагаемых, поэтому механически заменять слагаемые внутри произвольной суммы нельзя.
\par\addvspace{2pt}\begin{lrbox}{\CatFlowBox}\begin{minipage}[t]{\linewidth}
\setlength{\parskip}{2pt}\setlength{\columnsep}{4.5mm}\setlength{\intextsep}{0pt}
\begin{wrapfigure}[12]{r}{54mm}
\includegraphics[width=54mm,height=58mm,keepaspectratio]{images/unique/art-226.png}
\label{catflow:226}
\end{wrapfigure}
Переменное основание делит допустимую область на две части. В одной части основание больше единицы, в другой находится между нулём и единицей; точка основания, равного единице, всегда исключена. Прямой разбор этих частей и рационализация дают один результат, если их условия записаны полностью. Полезно сначала решить задачу двумя способами на простом примере и сравнить границы. После этого произведение, полученное рационализацией, читается обычным методом интервалов. Точки исчезнувших ограничений при этом продолжают отмечать отдельно. Даже если множитель сократился или стал квадратом, запрещённое основание не становится допустимым. Именно сохранение области определения превращает короткий приём в обоснованный метод решения.
\par\WFclear\end{minipage}\end{lrbox}%
\ifdim\dimexpr\ht\CatFlowBox+\dp\CatFlowBox\relax<59mm\dp\CatFlowBox=\dimexpr59mm-\ht\CatFlowBox\relax\fi%
\noindent\usebox{\CatFlowBox}\par\addvspace{2pt}
\begin{example}Решим \(\log_x(x+2)>2\). ОДЗ: \(x>0,x\ne1\). Рационализация даёт \((x-1)(x+2-x^2)>0\). Поскольку \(x+2-x^2=(2-x)(x+1)\), а \(x+1>0\), получаем \((x-1)(2-x)>0\). Ответ \((1;2)\). При прямом разборе \(0<x<1\) потребовалось бы \(x+2<x^2\), что невозможно; при \(x>1\) получаем тот же промежуток.\end{example}

\section{Системы и вложенные выражения}
\begin{example}Решим систему \(\log_2(x-1)\geq1\), \((\frac12)^{x-5}\geq1\). Первое условие равносильно \(x-1\geq2\), то есть \(x\geq3\). Во втором основание меньше единицы, поэтому \(x-5\leq0\). Пересечение даёт \([3;5]\), на котором все выражения определены.\end{example}
Для \(\log_2(\log_3(x-1))=1\) требуется \(x-1>0\) и \(\log_3(x-1)>0\); вместе это \(x>2\). Последовательно получаем \(\log_3(x-1)=2\), \(x-1=9\), \(x=10\). Систему логарифмических неравенств решают пересечением, совокупность — объединением; ни тот, ни другой переход не отменяет ОДЗ каждого исходного условия.

В системе с двумя неизвестными замена может вводиться для каждого из них. Если \(u=2^x,v=2^y\), то \(u,v>0\), и каждой допустимой паре соответствует ровно одна пара \((x,y)\).
\begin{example}Решим \(2^x+2^y=10\), \(2^{x+y}=16\). Для \(u=2^x,v=2^y\) получаем \(u+v=10,uv=16\). По Виету \(u,v\) являются корнями \(t^2-10t+16=0\), то есть \(2,8\) в любом порядке. Оба положительны; обратная подстановка даёт пары \((1,3)\), \((3,1)\).\end{example}

\relax

\par\addvspace{2pt}\begin{lrbox}{\CatFlowBox}\begin{minipage}[t]{\linewidth}
\setlength{\parskip}{2pt}\setlength{\columnsep}{4.5mm}\setlength{\intextsep}{0pt}
\begin{wrapfigure}[12]{r}{54mm}
\includegraphics[width=54mm,height=58mm,keepaspectratio]{images/unique/art-272.png}
\label{catflow:272}
\end{wrapfigure}
Сложное логарифмическое неравенство решают на общей области определения исходной записи. Положительность каждого аргумента, допустимость основания и ненулевой знаменатель проверяют до сокращения и замены переменной. Затем различают равенства выражений и сохранение их знаков: выражения одинакового знака можно заменять в произведении или частном для изучения знака, но нельзя произвольно подменять слагаемые суммы. Если новые критические точки появились после преобразования, их сопоставляют с исходной областью. В итоговом ответе открытые и закрытые концы имеют разное происхождение: равенство в неравенстве может быть разрешено, а точка с нулевым знаменателем или аргументом логарифма всё равно исключается.
\par\WFclear\end{minipage}\end{lrbox}%
\ifdim\dimexpr\ht\CatFlowBox+\dp\CatFlowBox\relax<58mm\dp\CatFlowBox=\dimexpr58mm-\ht\CatFlowBox\relax\fi%
\noindent\usebox{\CatFlowBox}\par\addvspace{2pt}
\relax
\section{Логарифмическое неравенство: вся цепочка ограничений}
\begin{example}[Полиномиальные аргументы и дробная правая часть]\label{ex:review-ege-exam-exponential-1}
Решим
\[
\log_3(x^2+4)-\log_3(x^2+x+1)
\ge\log_3\frac{x+3}{x+2}.
\]
Первый аргумент положителен при любом \(x\), второй равен \((x+\frac12)^2+\frac34>0\). Для правого аргумента требуется \(\frac{x+3}{x+2}>0\), то есть \(x<-3\) или \(x>-2\). Это исходная ОДЗ; точки \(-3,-2\) не могут вернуться в ответ после преобразований.
На ОДЗ объединим разность логарифмов и используем возрастание логарифма по основанию \(3\):
\[
\frac{x^2+4}{x^2+x+1}\ge\frac{x+3}{x+2},
\qquad
\frac{5-2x^2}{(x^2+x+1)(x+2)}\ge0.
\]
Положительный множитель \(x^2+x+1\) знак не меняет. Остались нули \(\pm\frac{\sqrt{10}}2\) и полюс \(-2\). Метод интервалов даёт \(( -\infty;-2)\cup[-\frac{\sqrt{10}}2;\frac{\sqrt{10}}2]\). Пересекаем именно с исходной ОДЗ:
\[
\boxed{\ x\in(-\infty;-3)\cup
[-\tfrac{\sqrt{10}}2;\tfrac{\sqrt{10}}2]\ }.
\]
Нули числителя включены, потому что знак нестрогий; точка \(-3\) исключена, потому что исходный логарифм нуля не существует.
\end{example}

\begin{example}[Сокращение не устраняет выколотую точку]\label{ex:review-ege-exam-exponential-2}
Решим
\[
\frac{(\log_2x)^2-4\log_2x+3}{\log_4x-\frac32}\ge0.
\]
ОДЗ: \(x>0\) и \(\log_2x\ne3\). При \(t=\log_2x\) получаем
\[
\frac{(t-1)(t-3)}{(t-3)/2}\ge0,
\qquad t\ne3.
\]
На этой области дробь равна \(2(t-1)\); следовательно, \(t\ge1\), \(t\ne3\). Обратный переход строго возрастает: \(x=2^t\). Ответ \([2;8)\cup(8;+\infty)\). Включение \(x=8\) после сокращения было бы ошибкой.
\end{example}
В составной задаче удобно вести два списка: ограничения исходной записи и ограничения конкретной замены. Замена \(t=a^x\) всегда требует \(t>0\); замена \(t=\log_a x\) при фиксированном допустимом основании принимает все действительные значения. Нулевой знаменатель исключается до сокращения, а при возврате к \(x\) учитывается направление монотонности.

\begin{example}[Показатель в числителе, многочлен в знаменателе]\label{ex:review-ege-exam-exponential-3}
Решим \(\frac{8^x-6\cdot4^x+12\cdot2^x-8}{(2x-3)^2}\ge0\). Знаменатель положителен при \(x\ne\frac32\). В числителе узнаём \((2^x-2)^3\). Его знак совпадает со знаком \(2^x-2\), значит, он неотрицателен при \(x\ge1\). Ответ \([1;\frac32)\cup(\frac32;+\infty)\). Замена \(t=2^x\) помогает разложить числитель, но заменять \((2x-3)^2\) на \((2t-3)^2\) нельзя: это другая функция. Нули и полюса обеих частей нужно совместить на одной оси \(x\).
\end{example}

\Needspace{7\baselineskip}\section*{Задачи}
\begingroup\setlength{\parskip}{1pt}
\noindent\begin{minipage}[t]{\linewidth}\textbf{23.1.} Реши \(5^{2x-1}=125\).\par\end{minipage}\par
\noindent\begin{minipage}[t]{\linewidth}\textbf{23.2.} Реши \((\dfrac14)^{x+1}=8\).\par\end{minipage}\par
\noindent\begin{minipage}[t]{\linewidth}\textbf{23.3.} Реши \(3^{x^2-4x}\geq1\).\par\end{minipage}\par
\noindent\begin{minipage}[t]{\linewidth}\textbf{23.4.} Реши \((\dfrac15)^{x^2-3x+2}<1\).\par\end{minipage}\par
\noindent\begin{minipage}[t]{\linewidth}\textbf{23.5.} Реши \(9^x-12\cdot3^x+27=0\).\par\end{minipage}\par
\par\addvspace{2pt}\begin{lrbox}{\CatFlowBox}\begin{minipage}[t]{\linewidth}
\setlength{\parskip}{2pt}\setlength{\columnsep}{4.5mm}\setlength{\intextsep}{0pt}
\begin{wrapfigure}[9]{r}{54mm}
\includegraphics[width=54mm,height=58mm,keepaspectratio]{images/unique/art-228.png}
\label{catflow:228}
\end{wrapfigure}
\textbf{23.6.} Реши \(4^x-6\cdot2^x+8\leq0\).\par

\textbf{23.7.} Реши \(2^{x+1}+2^x=24\).\par

\textbf{23.8.}\starmark{} Реши \(2^x+2^{-x}=\dfrac52\).\par

\textbf{23.9.} Реши \(3^{2x}-2\cdot3^x-3>0\).\par
\par\WFclear\end{minipage}\end{lrbox}%
\ifdim\dimexpr\ht\CatFlowBox+\dp\CatFlowBox\relax<58mm\dp\CatFlowBox=\dimexpr58mm-\ht\CatFlowBox\relax\fi%
\noindent\usebox{\CatFlowBox}\par\addvspace{2pt}

\noindent\begin{minipage}[t]{\linewidth}\textbf{23.10.} Реши \(\log_3(2x-1)=2\).\par\end{minipage}\par
\noindent\begin{minipage}[t]{\linewidth}\textbf{23.11.} Реши \(\log_{\dfrac12}(x+2)=-2\).\par\end{minipage}\par
\noindent\begin{minipage}[t]{\linewidth}\textbf{23.12.} Реши \(\log_2(x-1)+\log_2(x+1)=4\).\par\end{minipage}\par
\noindent\begin{minipage}[t]{\linewidth}\textbf{23.13.} Реши \(\log_3(x+4)-\log_3(x-2)=1\).\par\end{minipage}\par
\noindent\begin{minipage}[t]{\linewidth}\textbf{23.14.} Реши \((\log_2x)^2-3\log_2x+2=0\).\par\end{minipage}\par
\noindent\begin{minipage}[t]{\linewidth}\textbf{23.15.} Реши \(\log_{\dfrac12}(3-x)\geq-1\).\par\end{minipage}\par
\noindent\begin{minipage}[t]{\linewidth}\textbf{23.16.}\starmark{} Реши \(\log_2\dfrac{x-1}{x+1}\geq0\).\par\end{minipage}\par
\noindent\begin{minipage}[t]{\linewidth}\textbf{23.17.}\starmark{} Реши \(\log_2(x^2-4)\leq0\).\par\end{minipage}\par
\noindent\begin{minipage}[t]{\linewidth}\textbf{23.18.} Реши \(\log_x9=2\).\par\end{minipage}\par
\noindent\begin{minipage}[t]{\linewidth}\textbf{23.19.} Реши \(\log_{x-1}(x+3)=1\).\par\end{minipage}\par
\noindent\begin{minipage}[t]{\linewidth}\textbf{23.20.}\starmark{} Реши \(\log_x(x+2)>2\).\par\end{minipage}\par
\noindent\begin{minipage}[t]{\linewidth}\textbf{23.21.}\starmark{} Реши \(\log_{x+1}(x+3)\geq1\).\par\end{minipage}\par
\noindent\begin{minipage}[t]{\linewidth}\textbf{23.22.}\starmark{} Реши \(\log_{x+2}(x^2)>0\).\par\end{minipage}\par
\noindent\begin{minipage}[t]{\linewidth}\textbf{23.23.} Реши систему \(2^x\leq8\), \(\log_2(x-1)>0\).\par\end{minipage}\par
\noindent\begin{minipage}[t]{\linewidth}\textbf{23.24.} Реши систему \(\log_{\dfrac12}(x-1)\leq-1\), \(\log_2(5-x)\geq0\).\par\end{minipage}\par
\noindent\begin{minipage}[t]{\linewidth}\textbf{23.25.}\starmark{} Реши \(\log_2(\log_3(x-1))=1\).\par\end{minipage}\par
\noindent\begin{minipage}[t]{\linewidth}\textbf{23.26.} Реши \(2^{\log_2(x-1)}+\log_2 4=5\).\par\end{minipage}\par
\noindent\begin{minipage}[t]{\linewidth}\textbf{23.27.} Реши \(\log_2((x-1)^2)=2\).\par\end{minipage}\par
\noindent\begin{minipage}[t]{\linewidth}\textbf{23.28.}\starmark{} Реши \(\dfrac{\log_2(x-1)}{\log_2(x+1)}\geq0\).\par\end{minipage}\par
\noindent\begin{minipage}[t]{\linewidth}\textbf{23.29.}\starmark{} Реши \(\dfrac{\log_2x-1}{\log_2x+1}\geq0\).\par\end{minipage}\par
\noindent\begin{minipage}[t]{\linewidth}\textbf{23.30.}\starmark{} Реши систему \(2^x+2^y=12\), \(2^{x+y}=32\).\par\end{minipage}\par
\noindent\begin{minipage}[t]{\linewidth}\textbf{23.31.}\starmark{} Реши систему \(\log_2x+\log_2y=3\), \(x+y=6\).\par\end{minipage}\par
\noindent\begin{minipage}[t]{\linewidth}\textbf{23.32.}\starmark{} Реши \(\log_2(x^2+3)-\log_2(x^2+x+1)\ge\log_2\dfrac{x+2}{x+1}\), явно выписав исходную ОДЗ.\par\end{minipage}\par
\noindent\begin{minipage}[t]{\linewidth}\textbf{23.33.}\starmark{} Реши \(\dfrac{(\log_2x)^2-3\log_2x+2}{\log_4x-1}\ge0\). Объясни, какая точка остаётся запрещённой после сокращения.\par\end{minipage}\par
\noindent\begin{minipage}[t]{\linewidth}\textbf{23.34.}\starmark{} Реши \(\log_{x-1}(x^2-4x+5)\ge1\). Отдельно укажи ограничения на основание.\par\end{minipage}\par
\noindent\begin{minipage}[t]{\linewidth}\textbf{23.35.}\starmark{} Реши \(\dfrac{9^x-10\cdot3^x+9}{3^x-3}\le0\) через показательную замену.\par\end{minipage}\par
\noindent\begin{minipage}[t]{\linewidth}\textbf{23.36.}\starmark{} Реши \(\dfrac{\log_2(\log_3(x-1))}{\log_2(x-1)-1}\ge0\), сохранив условия существования вложенного логарифма.\par\end{minipage}\par
\noindent\begin{minipage}[t]{\linewidth}\textbf{23.37.}\starmark{} Реши \(\dfrac{27^x-12\cdot9^x+48\cdot3^x-64}{(3x-2)^2}\le0\). Объясни, почему знаменатель нельзя переписать как \((3t-2)^2\) при \(t=3^x\).\par\end{minipage}\par
\endgroup
